\section{Experiments}
\label{sec:experiments}
We compare answers and tool-workflow consistency while distinguishing configuration differences from attributable module effects. The illustrative case explains the mechanism and does not substitute for generalization or runtime-cost experiments.

\subsection{Setup and Main Results}
The main experiment covers all 603 GeoToolEval records; configuration comparisons use a fixed 57-record subset. FAS, TSF1, and TOS average per-task scores with failures retained. AgentGS uses GPT-5.6 Terra; main FAS values come from its first complete run, retaining partial credit. GeoAgent and EarthAgent use the same backbone. Direct-answer baselines are DeepSeek V4 Pro, GPT-5.6 Terra, GPT-5.4, Qwen-Plus, and GeoGPT. Process metrics compare successful calls from the three agents in the same tool environment against the same reference chains; direct-answer models have no tool traces.
\begin{table}[t]
\caption{GeoToolEval baseline comparison. AgentGS process metrics are per-record averages from the present 603-record end-to-end experiment; FAS values are retained from the source manuscript. Direct-answer language models have no tool traces; their process metrics are not applicable (---).}
\label{tab:main-results}
\centering
\begin{tabularx}{\columnwidth}{@{}Xrrr@{}}
\toprule
Method & FAS & TSF1 & TOS \\
\midrule
AgentGS (GPT-5.6 Terra) & \textbf{0.8731} & 0.9170 & 0.9166 \\
GeoAgent (GPT-5.6 Terra) & 0.7804 & 0.8663 & 0.7911 \\
EarthAgent (GPT-5.6 Terra) & 0.0862 & 0.1182 & 0.1390 \\
DeepSeek V4 Pro & 0.5162 & --- & --- \\
GPT-5.6 Terra & 0.3137 & --- & --- \\
GPT-5.4 & 0.3069 & --- & --- \\
Qwen-Plus & 0.2918 & --- & --- \\
GeoGPT & 0.2500 & --- & --- \\
\bottomrule
\end{tabularx}
\end{table}


AgentGS obtains FAS of 0.8731 versus GeoAgent's 0.7804 and direct GPT-5.6 Terra's 0.3137 (Table~\ref{tab:main-results}). Its present end-to-end run over 603 records yields TSF1/TOS of 0.9170/\allowbreak 0.9166, including execution failures and incomplete final answers. The table reports 0.8663/\allowbreak 0.7911 for GeoAgent and 0.1182/\allowbreak 0.1390 for EarthAgent. These scores measure reference-tool coverage and relative order; correct inputs, artifacts, and scientific results require additional checks.

Single-step, two- to three-step, and at-least-four-step FAS values are 0.9425, 0.7557, and 0.6167. On 119 two- to three-step records, TSF1/TOS are 0.9779/\allowbreak 0.9832; on 60 longer records, they are 0.8013/\allowbreak 0.7785. Lower long-chain process scores indicate remaining coverage and ordering difficulties. Longer chains are concentrated in remote sensing, so domain and chain-length effects overlap (Figure~\ref{fig:performance}).
\begin{figure*}[t]
  \centering
  \includegraphics[width=\textwidth]{performance.pdf}
  \caption{Overall and domain-specific performance of AgentGS and baselines on GeoToolEval: (a) mean FAS across all evaluation records; (b) mean FAS across the ten geoscientific domains.}
  \label{fig:performance}
  \Description{A horizontal bar chart compares mean FAS for AgentGS, two agent baselines, and five language models. A radar chart presents performance across ten geoscientific domains.}
\end{figure*}


\subsection{External Evaluation on GeoBenchX}
\label{sec:external-evaluation}
We integrate official GeoBenchX tools~\cite{krechetova2025geobenchx} and construct dedicated Skills, Capabilities, dependencies, and mappings. AgentGS and the official ReAct baseline are each reported on 202 tasks using GPT-5.6 Terra. That model is also the single judge: 0, 1, and 2 denote mismatch, partial match, and match to reference solutions. Mean score and full-match rate are separate from GeoToolEval FAS.
\begin{table}[t]
\caption{Evaluation results on GeoBenchX. Both methods use GPT-5.6 Terra.}
\label{tab:external-results}
\centering
\begin{tabularx}{\columnwidth}{@{}Xrr@{}}
\toprule
Method & Mean score & Full-match \\
 & $(0$--$2)\uparrow$ & rate $(\%)\uparrow$ \\
\midrule
ReAct & 1.1980 & 53.47 \\
AgentGS & \textbf{1.3762} & \textbf{62.87} \\
\bottomrule
\end{tabularx}
\end{table}


AgentGS scores 1.3762 versus 1.1980 for ReAct, a difference of 0.1782 points; full-match rates are 62.87\% versus 53.47\%, a difference of 9.40 percentage points. Although task counts are supplied, exact task membership and consistent experimental conditions still require verification before interpreting these differences as transfer evidence.

\subsection{Ablation Configuration Comparisons}
\label{sec:ablations}
The subset covers all ten domains, with 22 single-step, 20 two- to three-step, and 15 longer records; multistep tasks comprise 61.40\%. The design uses Full on the same subset as reference, but actual intervention scopes are limited: A1 is a partial control, A2's intended ablation was not activated, and A5 changes only part of ranking. Differences therefore do not uniformly represent controlled single-module effects.

A1 is designed to plan at the tool level with the same strategies, dependencies, and tools. A2 is designed to disable mandatory-predecessor expansion while retaining selection and dependency declarations; completion may compensate for omissions. A3 retains gap diagnosis but disables supporting-capability insertion, including bridging. A4 disables only bridging. A5 is designed to choose the first valid tool in a fixed order while retaining requirements, scopes, and validation. A6 disables local recovery, retaining planning, grounding, and binding before the first attempt.
\begin{table}[t]
\caption{Earlier component comparison on 57 records, including 35 multistep records. A1 is a partial control; A2 did not activate its intended intervention.}
\label{tab:ablation}
\centering
\begin{tabularx}{\columnwidth}{@{}Xrr@{}}
\toprule
Configuration & FAS & Multistep FAS \\
\midrule
Full & 0.6387 & 0.5333 \\
A1: Flat planning & 0.4748 & 0.4095 \\
A2: w/o predecessor expansion & 0.4760 & 0.4095 \\
A3: w/o automatic completion & 0.4624 & 0.3802 \\
A4: w/o directed bridging & 0.6160 & 0.4913 \\
\bottomrule
\end{tabularx}
\end{table}

\begin{table*}[t]
\caption{Ablation process metrics. Multistep scores cover the 35 subset records with at least two reference calls. Each column averages per-task scores over its corresponding records. Full uses the supplementary 57-task experiment; A1--A6 use per-task records from the formal 57-task worksheet.}
\label{tab:ablation-process}
\centering
\begin{tabularx}{\textwidth}{@{}Xrrrr@{}}
\toprule
Configuration & Mean TSF1 & Mean TOS & Multistep TSF1 & Multistep TOS \\
\midrule
Full method & 0.9249 & 0.9392 & 0.8862 & 0.8952 \\
A1: Direct tool-level planning & 0.8599 & 0.8778 & 0.8147 & 0.8295 \\
A2: No predecessor expansion & 0.8694 & 0.8958 & 0.8302 & 0.8588 \\
A3: No independent completion & 0.8732 & 0.8865 & 0.8364 & 0.8438 \\
A4: No directed bridging & 0.8603 & 0.8997 & 0.8249 & 0.8652 \\
A5: No tool-selection optimization & 0.8503 & 0.8851 & 0.8133 & 0.8414 \\
A6: No execution recovery & 0.8208 & 0.8588 & 0.7652 & 0.7985 \\
\bottomrule
\end{tabularx}
\end{table*}


A1's FAS decreases from Full are 16.39/\allowbreak 12.38 percentage points; A2's are 16.27/\allowbreak 12.38 and A3's are 17.63/\allowbreak 15.31 (Table~\ref{tab:ablation}). Multistep decreases are smaller in all three cases. The partial A1 control and inactive A2 ablation prevent direct attribution to hierarchical planning or predecessor expansion. The comparison does not show greater multistep effects for predecessor expansion or independent completion.

A4's overall/multistep decreases are 2.27/\allowbreak 4.20 points; A6's are 4.15/\allowbreak 7.20. Attribution to bridging or recovery requires inserted branches, failure/recovery traces, and consistent conditions. A5 increases overall FAS from 0.6387 to 0.6415 but lowers multistep FAS from 0.5333 to 0.5184, giving different directions for its partial ranking change.

Full TSF1/TOS are 0.9249/\allowbreak 0.9392 overall and 0.8862/\allowbreak 0.8952 on 35 multistep records (Table~\ref{tab:ablation-process}). Full uses the supplementary experiment; A1--A6 use formal worksheet records. All six have lower overall and multistep process scores. Multistep changes from Full are $-7.15$/\allowbreak $-6.57$ points for A1, $-5.60$/\allowbreak $-3.64$ for A2, and $-4.98$/\allowbreak $-5.14$ for A3. A4 changes by $-6.13$/\allowbreak $-3.00$ and has lower TSF1 but higher TOS than A3. A5 changes by $-7.29$/\allowbreak $-5.38$, reflecting partial ranking adjustment. A6 changes by $-12.10$/\allowbreak $-9.67$, with a larger coverage decrease. These are configuration differences; parameters, artifact consumption, and module effects require traces and scrutiny of data provenance.

\subsection{Illustrative Case Study}
\label{sec:case-study}
Consider several rasters of the same region and variable, with aligned grids and a reference value in the same units. The task requires a pixel-wise mean composite, its valid-pixel mean, the difference from the reference, and the comparison relation.

Skill Retrieval returns RasterCompose ($A$) and ReferenceCompare ($B$). Planning selects InputCheck ($w$) and RasterMean ($u$) from $A$, and ScalarDiff ($v$) and SignCheck ($z$) from $B$, yielding $w\to u$ and $v\to z$. Composition expands required predecessors while retaining provenance and tool scopes. The fragments remain disconnected: $u$ produces a raster file, whereas $v$ requires a numerical value.

\begin{figure*}[t]
\centering
\includegraphics[width=0.64\textwidth]{case_study.pdf}
\caption{Illustrative cross-Skill completion. Selected fragments $w\to u$ and $v\to z$ are connected through inactive Skill $C$'s $u\to m\to v$ branch. The resulting five-step workflow is grounded in tools with explicit artifact propagation.}
\label{fig:case-study}
\Description{A mean-composite task is planned using two Skills, completed using valid-pixel mean statistics from a third Skill, and executed as input validation, compositing, statistics, differencing, and sign comparison.}
\end{figure*}

Completion retrieves $u\to m\to v$ from inactive RasterDiff ($C$), where ValidMean ($m$) computes the composite's valid-pixel mean. The query explicitly requires this statistic. Under the example's assumptions, this is the unique branch satisfying prerequisites, types, availability, and acyclicity. Insertion retains $u$'s constraints, adds $m$, updates $v$'s dependencies, tools, and bindings, and preserves $w\to u$ and $v\to z$, forming $w\to u\to m\to v\to z$ (Figure~\ref{fig:case-study}).

Grounding selects permitted validation, compositing, statistics, differencing, and comparison tools. Bridge grounding requires unique endpoint instances with the correct order. Execution passes the validated raster list to compositing, then passes a list containing only the composite path to statistics. The returned mean and task reference bind to $a$ and $b$ for $d=a-b$; sign comparison evaluates $d$ against zero. Thus, each output concerns the same composite and reference.

If statistics fails, local recovery retains validation and compositing and reuses the artifact. Uncertain calls are not replayed, and recovery requiring another capability suspends execution. Successful results yield the raster, mean, difference, and relation. After verification, the workflow can form candidate CompositeCompare ($D$), with rebindable image-list/reference inputs, artifact dependencies, and usage constraints. Its applicability and reuse benefit on other tasks require further validation.
